\begin{abstract}
Vibration-based fault diagnosis constitutes a fundamental aspect of predictive maintenance in industrial systems. However, conventional approaches encounter substantial difficulties arising from data scarcity, domain discrepancies, and constrained generalization across diverse machinery types. This study proposes VibSemAlign, an innovative framework that employs multimodal large language models (MLLMs) to integrate vibration signals with semantic descriptions of faults, thereby facilitating robust zero-shot and few-shot diagnostics. Raw vibration data are transformed into spectrograms, which are then associated with textual prompts delineating fault categories. A specialized vibration-semantic contrastive loss refines vision-language models, establishing strong correspondences between visual signal features and natural language semantics. When combined with model-agnostic meta-learning (MAML), the framework supports efficient adaptation to novel domains using limited samples. Comprehensive evaluations on CWRU and Paderborn benchmarks reveal pronounced advantages, including accuracy increases of up to 15\% in cross-domain zero-shot contexts relative to state-of-the-art deep learning methods. VibSemAlign surpasses traditional signal processing techniques and CNN-based classifiers, while also delivering interpretable semantic explanations for detected faults. The primary contributions comprise the initial MLLM-based semantic alignment for vibration analysis, a customized contrastive objective, and empirical confirmation of generalization across industrial datasets. These advancements lay the groundwork for semantically informed, data-efficient fault diagnosis in operational environments.
\end{abstract}